\documentclass[12pt]{report}

\usepackage{blindtext}
%\usepackage[T1]{fontenc}
\usepackage[francais]{babel}
\usepackage[utf8]{inputenc}
%\usepackage[pdftex]{graphicx}
%\usepackage{xcolor}
\usepackage{enumitem}
%\usepackage{listings}
\usepackage{graphicx}
%\usepackage[export]{adjustbox}
%\usepackage{wrapfig}
\usepackage{titlesec}
%\usepackage{hyperref}
%\usepackage{ulem}


\usepackage[a4paper,left=2cm,right=2cm,top=2cm,bottom=2cm]{geometry}
\usepackage{libertine}

\setlength{\parindent}{0cm}
\setlength{\parskip}{1ex plus 0.5ex minus 0.2ex}
%\newcommand{\hsp}{\hspace{20pt}}
\newcommand{\HRule}{\rule{\linewidth}{0.5mm}}


\titleformat{\chapter}[display]
  {\normalfont\bfseries}{}{0pt}{\Large}


\begin{document}

\begin{titlepage}
  \begin{sffamily}
  \begin{center}
    \begin{minipage}{1\textwidth}
    \end{minipage}
    \newline
    \newline
    \newline
    \newline
    \newline
    \newline
    \newline
    \newline
    \newline
    \newline
    \newline
    \textsc{\LARGE Université de Valenciennes et du Hainaut-Cambrésis}\\[4cm]
    \HRule \\[0.4cm]
    { \huge \bfseries M2C - Malware Clustering et Classification\\[0.4cm] }
    \HRule \\[2cm]
    \begin{figure}[htpb]
    \center
	\includegraphics[width=16cm,height=8cm]{img/datascience.png}
    \end{figure}
     \center
        Vincent \textsc{Romé} -
        Axel \textsc{Foulon} -
        Julien \textsc{Dupagny}
    \vfill
    {\large — 9 Juin 2018 —}
  \end{center}
  \end{sffamily}
\end{titlepage}
\newpage
\begin{center}
\begin{tabular}{|l|c|l|c|}
   \hline
   \multicolumn{4}{|c|}{HISTORIQUE DES VERSIONS} \\
   \hline
   DATE & VERSION & ÉVOLUTION DU DOCUMENT & RÉDACTEUR \\
   \hline
   9/06/2017 & 0.1 & Version préliminaire & Équipe complète \\
   \hline
\end{tabular}
\end{center}
\tableofcontents
\chapter{Introduction}
	Jusqu'à présent les systèmes de détection d'intrusion reposaient traditionnellement sur des signatures générées manuellement par des experts en sécurité, puis nous avons vu apparaître des systèmes permettant de détecter des patterns entre des jeux de données ceci à permis de générer automatiquement ces règles. Aujourd’hui avec l’apparition du “Big Data” des solutions comme le Deep Learning ou  Machine learning sont souvent présentés comme les technologies pouvant révolutionner les systèmes de détections et les performances.
Ces systèmes permettent de générer automatique un modèle de détection à partir de données et leurs capacité à généraliser les événements malveillants permettait en effet de détecter des éléments encore inconnus.
	L’objectif est d’ici comprendre le fonctionnement de ces algorithmes et de les appliquer au milieu de la sécurité informatique. D’exposer les résultats de notre recherche sur la détection de fichier PE (exécutable windows malveillants). Tout en gardant un regard critique sur les résultats et en essayant d’apporter des solution sur l’utilisation de tel algorithme en production.
 
Un modèle de détection supervisé est construit à partir de données labellisées fournies par l’expert : des événements bénins, mais aussi des événements malveillants pour guider le modèle de détection. L’algorithme d’apprentissage va automatiquement chercher les points permettant de caractériser chacune des classes ou de les discriminer pour construire le modèle de détection


\chapter{Abstract}
\chapter{Contexte du projet}
\end{document}